\subsection{Odd internal roots of minimum one}
\label{proof:Q5-C186}

Take ordered root lengths
\[
 3\le a<b<s_1<t_1<\cdots<s_k<t_k,\qquad k\ge0,
 \tag{Q186.0.2}
\]
and a terminal residual list $\Sigma$ of at least two distinct
lengths. Orient its arms by Table~\ref{tab:odd186:directions}.
For genuine equal-pair decorations $E_o,E_t$, set
\[
 H_k=E_o pW_a\bar f_b
          \prod_{i=1}^kW_{s_i}\bar f_{t_i}\,
          T^L(E_tF_{\Sigma,L}).
 \tag{Q186.0.3}
\]
All means and operators use the original finite probability measure
and the original skew kernel; $d=\E u^2$.

\begin{proposition}\label{prop:odd-complete:minimum-one-root}
For every such residual list, every finite $k$ and all independent
positive integers $L,M$,
\[
 \E H_0^4\le e^{-27d/5},\qquad
 \tag{Q186.0.4}
\]
\[
 \E H_k^4\le e^{-[27/5+\zeta_k]d}\quad(k\ge1),\qquad
 \zeta_k=2737/3840+64(k-1)/5.
 \tag{Q186.0.5}
\]
Writing $\lambda_0=27/5$ and $\lambda_k=27/5+\zeta_k$ for
$k\ge1$, the outgoing message satisfies
\[
 I_q(T^MH_k)^3\le e^{-[2\lambda_k-c(M)]d}\le e^{-6d/95},
 \qquad I_q(X)=\E[q^{-1/3}|X|^{8/3}],
 \tag{Q186.0.6}
\]
where $c(1)=204/19$, $c(2)=10$, $c(3)=0$, $c(4)=7$ and
$c(M)=6$ for $M\ge5$.
\end{proposition}

\paragraph{Q186.1 Literal directions and exhaustive terminal table}\label{part:Q186.1}

An \(f_n\) arm has every edge away from its root; \(\bar f_n\) has every edge toward its
root; \(h_n\) has its first edge toward the root and its remaining edges away
from the first neighbor; \(W_n\) has its first edge away from the root and its
remaining edges toward that neighbor. \(T^L\) means every original inner edge
points from the terminal toward the internal root. \(T^M\) means every outer
edge points from that root toward the receiving center. The original unit
arm at the internal root points away from the root, giving p.

In the table, every displayed residual list is strictly increasing. The
indices used for terminal lengths are independent of the root a, b.

% Replace only q5-c186.tex's description at baseline lines 81--101.
% Keep the surrounding part:Q186.1 anchor and its explanatory text.
\begingroup\small
\renewcommand{\arraystretch}{1.15}
\begin{longtable}{@{}>{\raggedright\arraybackslash}p{.44\linewidth}
                      >{\raggedright\arraybackslash}p{.50\linewidth}@{}}
\caption{Terminal directions at a minimum-one internal root.
All residual lists are strictly increasing; their indices are independent
of the internal root lengths \(3\le a<b\).
The root pair is always \(W_a\bar f_b\).}
\label{tab:odd186:directions}\\
\toprule
Terminal residual and length condition&Actual terminal message \(F\)\\
\midrule\endfirsthead
\toprule
Terminal residual and length condition&Actual terminal message \(F\)\\
\midrule\endhead
\bottomrule\endfoot
\((1,t)\), \(t\ge2\)&\(qf_t\)\\
\((2,t)\), \(t\ge4\)&\(\bar f_2W_t\)\\
\((2,3)\)&\(h_2f_3\)\\
\((3,t)\), \(t\ge5\), \(L=1\)&\(h_3f_t\)\\
\((3,t)\), \(t\ge5\), \(L\ge2\)&\(\bar f_3W_t\)\\
\((3,4)\)&\(h_3f_4\)\\
Even \((\alpha_1,\beta_1,\ldots,\alpha_j,\beta_j)\),
\(\alpha_1\ge4\), \(j\ge1\)
  &\(\prod_{i=1}^j h_{\alpha_i}f_{\beta_i}\)\\
The same even list, \(\alpha_1\le3\), \(j\ge2\)
  &\(\prod_{i=1}^j h_{\alpha_i}f_{\beta_i}\)\\
\((1,s,t)\), \(1<s<t\)&\(qf_sh_t\)\\
\((1,\alpha_1,\beta_1,\ldots,\alpha_j,\beta_j)\), \(j\ge2\),
\((\alpha_1,\beta_1,\alpha_2)=(2,3,4)\)
  &\(\begin{aligned}[t]
       &qf_2h_{\beta_2}h_3f_4\\
       &\quad\cdot\prod_{i=3}^j h_{\alpha_i}f_{\beta_i}
     \end{aligned}\)\\
The same minimum-one odd list, all other cases
  &\(qf_{\alpha_1}h_{\beta_1}
        \prod_{i=2}^j h_{\alpha_i}f_{\beta_i}\)\\
Nonunit odd \((s,\beta_1,\gamma_1,\ldots,\beta_j,\gamma_j)\),
\(s\ge2\), \(\beta_1=3\) (hence \(s=2\))
  &\(\bar f_2 f_3h_{\gamma_1}
        \prod_{i=2}^j h_{\beta_i}f_{\gamma_i}\)\\
The same nonunit odd list, \(\beta_1\ge4\)
  &\(\bar f_s\prod_{i=1}^j h_{\beta_i}f_{\gamma_i}\)\\
\end{longtable}
\endgroup


The fixed cross in the tenth row is part of the original direction, not an
analytical switch. Every list with at least two residual arms occurs in exactly
one row after the displayed conditions are applied. The separate (3, t) rows
depend only on whether \(L=1\). No direction depends on d, \(\mu\), K or a frequency.
For all of them the new internal root direction is exactly \(W_a \bar f_b\).

\begin{proof}
\paragraph{Complete root and terminal certificates}
\label{part:Q186.3}\label{part:Q186.4}
Use the full constant split and the minimum-one root transfer in
Section~\ref{sec:odd-complete:comparison}. The positive finite
propagation in each terminal family provides its own $C,c,Q$;
the transfer is applied to those complete parameters.

\paragraph{Q186.5 Two-arm terminals $(1,t)$}\label{part:Q186.5}
For every $3\le a<b$, $t\ge2$ and $L\ge1$, the actual message is
$H_0=pW_a\bar f_bT^L(qf_t)$. Apply
Proposition~\ref{prop:root-post:square-moments} with one root pair
$(u_1,v_1)=(a,b)$ and terminal parameter $t$. It gives
\fitdisplay{
 \E H_0^4\le e^{-705d/128}.
}{Q186.5.1}
The root and terminal lengths are independent in that proposition.
Its pointwise conclusion \eqref{eq:root-post:canonical-certificate}
also supplies the whole vector required for later root pairs. In the
present $B=-K$ notation, set $V=-Z$. Then
\begin{equation}\label{eq:odd186:short-one-certificate}
 \begin{gathered}
 H_0\le e^{a_0+BV},\qquad \|V\|^2\le9d/4,\\
 4a_0+4\E BV+(17/4)\|V\|^2\le-705d/128.
 \end{gathered}
\end{equation}
There is no separate curvature term in this certificate. It holds
on the entire host domain, including all short inner lengths, and
therefore applies to the low-domain restoration in part~\ref{part:Q186.11}.
For $a\ge4$, the direct fourth payment in
part~\ref{part:root-post:payments} also gives the stronger loss $7d$.

\paragraph{Q186.6 All remaining two-arm terminals of minimum two or three}

\paragraph{Q186.6.1 Every terminal \((2,t)\), \(t\ge4\)}\label{part:Q186.6.1}
Use the same actual root \(pW_a\bar f_b\) and terminal
\(F=\bar f_2W_t\). The minimum-one instance of
Theorem~\ref{thm:odd-two:roots} proves, for all \(3\le a<b\)
and every inner length,
\begin{equation}\label{eq:odd186:terminal-two}
 \E[pW_a\bar f_bT^LF]^4\le e^{-6d}.
\end{equation}
Its complete vectors have norm at most \(2\sqrt d\), with any
used \(p\)-curvature payment kept separate. Thus this is also
an input for the paid restoration of later root pairs below.

\paragraph{Q186.6.2 The actual terminal \((2,3)\)}
Choose \(F=h_2f_3\), as in the direction table. Its actual arm
sum is \(2-d-B^2u\), equal to the arm sum for \(\bar f_2W_3\).
Theorem~\ref{thm:odd-two:roots} treats this genuine alternative
with the same \(6d\) fourth loss and the same complete-vector
bound. It uses this common positive envelope, not an equality of
the two different terminal products.

\paragraph{The remaining terminal families}
\label{part:Q186.7}\label{part:Q186.8}\label{part:Q186.9}
Theorem~\ref{thm:odd-three:roots} proves loss six for the actual
$(3,t)$ directions, including the fixed $(3,4)$ direction.
Proposition~\ref{prop:odd-even:roots} proves loss eight for even
lists of minimum at least four and at least $11/2$ for the other
even multipair lists. The $q$-first-row packets of
Section~\ref{sec:odd-complete:q-terminals} give $705/128$ for
every minimum-one odd terminal, with the fixed cross prescribed
in the table. Section~\ref{sec:odd-complete:inward-anchor} gives
$13/2$ for every nonunit-minimum odd terminal. All these are
complete fourth certificates at the same root $pW_a\bar f_b$.

\paragraph{Complete ordered root and outgoing connector}
\label{part:Q186.11}
The losses in \eqref{eq:odd-complete:catalogue} are at least
$27/5$, proving \textup{(Q186.0.4)}.
Lemma~\ref{lem:odd-complete:vector-cap} supplies the pointwise
low-domain certificate with its full vector norm at most
$(7/2)\sqrt d$. The paid restoration in
Section~\ref{sec:odd-complete:assembly} gives the additional
$\zeta_k$ and restores the genuine equal pairs.
Finally \eqref{eq:odd-common:outlet} applies to the actual real
message with the independent outgoing length $M$.
Since $54/5-204/19=6/95$, it gives \textup{(Q186.0.6)}.
\end{proof}

